\documentclass[11pt]{article}
\usepackage[hcentering=true, width=6.75in, height=9.5in, paper=letterpaper]{geometry}
\usepackage{lineno,hyperref}
\usepackage[inkscapelatex=false]{svg}
\usepackage{multirow}
\usepackage{amsmath}
\usepackage{amssymb}
\numberwithin{equation}{section}
\usepackage{graphicx}
\usepackage{mathtools}
\usepackage{float} 
\usepackage{tikz}
\usetikzlibrary{decorations.pathreplacing, arrows.meta, positioning, backgrounds, fit,calc}
\usepackage{enumerate}
\usepackage{booktabs}
\usepackage{amsfonts}
\usepackage{bm}
\usepackage{caption}
\usepackage{subcaption}
\usepackage{comment}
\usepackage{diagbox}
\usepackage{rotating} 
\usepackage{bbm}
\usepackage{algorithm}
\usepackage{algpseudocode}
\usepackage{xcolor}
\usepackage{adjustbox}
\usepackage{threeparttable}
\usepackage[table]{colortbl}
\usepackage{ulem}
\usepackage{natbib}
\usepackage{amsthm,setspace}
\usepackage{needspace}

\modulolinenumbers[5]
\newtheorem{myDef}{Definition} 
\newtheorem{myTheo}{Theorem}
\newtheorem{myExam}{Example}
\newtheorem{lemma}{Lemma}
\newtheorem{proposition}{Proposition}
\newtheorem{corollary}{Corollary}
\newtheorem{assum}{Assumption}
\newtheorem{remark}{Remark}
% Locate shared bibliography files from either the project root or backup/.
\IfFileExists{./model5-names.bst}{%
  \bibliographystyle{model5-names}%
  \newcommand{\bibliographyfile}{references}%
}{%
  \bibliographystyle{../model5-names}%
  \newcommand{\bibliographyfile}{../references}%
}
\doublespacing

\begin{document}
\newcommand{\argmin}{\operatornamewithlimits{argmin}}
\newcommand{\mc}[1]{\mathcal{#1}}

\title{Neural Regression Trees for Censored and Uncensored Outcomes}
\author{Yixing Zhao, Liqun Diao}
\maketitle
\begin{abstract}
Estimation of the full conditional distribution is considered for exact or censored continuous outcomes. Nonparametric regression allows this distribution to vary with covariates without imposing a fixed parametric form. A Neural Regression Tree (NRT) is proposed for this purpose. The response is partitioned by a complete binary tree with a fixed topology on a bounded time window, with an additional terminal node for the probability beyond that window. For each covariate vector, body split locations and left-branch probabilities are determined by a neural network; the final cut at the horizon is fixed. Products of the branching probabilities along each path give normalized terminal-node masses. A uniform density within each terminal body interval, together with the tail mass, defines a valid conditional law and distribution and survival functions on the analysis window. The validity of this construction is established. As the tree and network complexities increase, continuous conditional distribution functions on compact domains can be approximated uniformly under the stated assumptions. Simulation studies and real-data analyses are used to evaluate estimation and prediction performance in uncensored and mixed-censored settings.
\end{abstract}

\section{Introduction}

Regression methods are commonly used to estimate how the conditional mean of a response varies with covariates. However, the mean alone does not describe tail behavior, multimodality, or subject-specific risk profiles. These features may vary with covariates in ways that cannot be summarized by a single location parameter. Conditional density estimation and distributional regression address this limitation by modeling the full conditional distribution of the response.

Survival regression models the conditional distribution of event times given covariates. Estimation is often complicated by censoring, and different subjects may be observed under different schemes, including right censoring, interval censoring, and current-status observation \citep{lawless_statistical_2003,sun_statistical_2006, kleinbaum_survival_2012}. Traditional survival models, such as the accelerated failure time (AFT) model \citep{lawless_statistical_2003}, often require parametric assumptions on the residual distribution, whereas the Cox proportional hazards model \citep{cox_regression_1972} relies on the proportional hazards assumption. These models are widely used because they are interpretable and computationally efficient. When their assumptions are violated, estimates may be biased and prediction may be poor.

Tree methods have been developed for distribution fitting. In these methods, the sample space is recursively partitioned, and probability mass is allocated across nested regions to describe broad features at coarse levels and local features at finer levels. P\'olya tree priors provide a Bayesian construction in which beta distributions are assigned to the branching probabilities on a fixed recursive partition \citep{mauldin_polya_1992,lavine_aspects_1992}. Tree-based regression models describe complex covariate effects through a sequence of local decisions. In a standard regression tree, the covariate space is recursively partitioned, and a prediction, usually a node-specific mean, is assigned to each terminal node. Splits are often chosen to reduce squared-error loss \citep{taylor_splitting_1996,loh_classification_2011}. Distributional trees extend this construction by estimating a response distribution within each terminal node. The node-specific model may be a parametric family whose location, scale, and shape vary across covariate regions, or a nonparametric model such as a histogram \citep{zeileis_model-based_2008,schlosser_distributional_2019,yang_conditional_2024}. Tree-based estimates can be sensitive to the choice of partition. Bayesian priors on tree structures were considered by \citet{wong_optional_2010,awaya_hidden_2024} to allow the partition to adapt to the data. These methods were developed for density estimation, and additional modeling is needed to make the response distribution depend on covariates. Forest methods average many randomized trees to improve stability and prediction and can be viewed as adaptive local-likelihood estimators of the conditional distribution \citep{schlosser_distributional_2019,hothorn_predictive_2021}. Tree-based methods have also been developed for censored outcomes, including random survival forests \citep{ishwaran_random_2008} and tree or forest methods for interval-censored data \citep{fu_survival_2017,yao_ensemble_2021}. More recently, \citet{angelim_partition_2026} construct data-adaptive partitions of the joint covariate-outcome space and select splits using conditional negative log-likelihood.

Neural networks have been used to model complex covariate effects in survival analysis \citep{sun_neural_2023,liu_neural_2025,gao_neural_2025,yuan_interpretable_2026}. In these models, multilayer nonlinear transformations are used to represent nonlinear effects and high-order interactions without prespecifying their functional form. Several methods have been developed for censored data. In \cite{katzman_deepsurv_2018}, a deep neural network is used to model nonlinear covariate effects and interactions within a proportional hazards model. For interval-censored data, a neural-network model was developed to predict Alzheimer's disease progression after genetic markers were screened under interval censoring and left truncation \cite{sun_neural_2023}. A neural interval-censored regression model with feature selection was proposed for high-dimensional predictors in \cite{meixide_neural_2024}. Neural networks have also been combined with dynamic prediction and partially linear transformation models to accommodate longitudinal covariates and nonlinear nuisance effects while retaining interpretable model components \citep{liu_neural_2025,yuan_interpretable_2026}. Applications outside biomedicine include prediction of remaining useful life from censored degradation signals \citep{gao_neural_2025}. These methods support the use of neural networks for censored survival prediction. However, many retain a specific hazard structure, transformation model, or parametric lifetime family. The conditional distribution is therefore still restricted by the chosen model, and many methods do not provide a common probabilistic formulation for mixed censoring schemes.

Neural networks have also been combined with partitions of the response space. \citet{li_deep_2021} formulate conditional distribution estimation as ordered classification over fixed response bins and introduce a joint binary cross-entropy loss. \citet{hu_conditional_2023} estimate a conditional hazard using a neural network and a full-likelihood loss for uncensored and right-censored outcomes. These methods use fixed response grids for neural distribution regression. In this paper, we propose a Neural Regression Tree (NRT) for conditional distribution estimation with exact or censored continuous outcomes. A tree is used to partition the response space recursively, with both cut locations and branching probabilities determined by the covariates. Products of local branching probabilities define the terminal masses of a conditional distribution. Exact observations contribute a density to the likelihood, whereas right-censored, interval-censored, and current-status observations contribute probabilities assigned to their observed response sets.

The remainder of the paper is organized as follows. The NRT construction, neural parameterization, and penalized likelihood estimation are described in Section~\ref{sec:NRT}. Section~\ref{sec:theory} studies the validity, approximation, identification, and consistency. Simulation studies are presented in Section~\ref{sec:simulation}, followed by real-data analysis in Section~\ref{sec:rwd} and discussion in Section~\ref{sec:discuss}. 

\section{Neural Regression Tree}
\label{sec:NRT}
The neural regression tree is constructed by combining a recursive response partition with covariate-dependent branching probabilities.

\subsection{Regression Tree}
\label{section-regression-tree}
Let $T\in\mathcal{S}\subseteq\mathbb{R}^+$ be a response and let $X\in\mathcal{X}\subseteq\mathbb{R}^d$ be a vector of covariates. Consider a recursive binary partition of the response space that may depend on the covariates,
$$
    \Pi(x)=\bigcup_{m\geq0}\mathcal{N}_m(x),
    \qquad
    \mathcal{N}_m(x)=\{A_{m,1}(x),\ldots,A_{m,2^m}(x)\},
$$
where $\mathcal{N}_0(x)=\{A_{0,1}(x)\}$ with $A_{0,1}(x)=\mathcal{S}$. For each $m\geq0$, the sets in $\mathcal{N}_m(x)$ form a disjoint partition of $\mathcal{S}$, and each parent node is split into two children,
$$
    A_{m,j}(x)=A_{m+1,2j-1}(x)\cup A_{m+1,2j}(x),
    \qquad
    A_{m+1,2j-1}(x)\cap A_{m+1,2j}(x)=\emptyset .
$$
Thus, the tree is defined on the range of the response, but the cut locations and hence the recursive partition may vary with the covariates.

The binary topology fixes the node labels and their parent-child relations across covariate values. A node label $A$ is distinguished from its interval $A(x)=(a_A(x),b_A(x)]$. Write $W_A(x)=b_A(x)-a_A(x)$ for the interval width. For an internal node, let $\rho_A(x)\in(0,1)$ denote the cut proportion. The split point is
$$
    c_A(x)=a_A(x)+\rho_A(x)W_A(x),
$$
with left and right children
$$
    A^\ell(x)=(a_A(x),c_A(x)],
    \qquad
    A^r(x)=(c_A(x),b_A(x)].
$$
The covariates also enter through the conditional branching probability from the current node to its left child,
$$
    \theta_A(x)
    =
    P\{T\in A^\ell(x)\mid T\in A(x),X=x\},
    \qquad 0<\theta_A(x)<1,
$$
so that
$$
    P\{T\in A^r(x)\mid T\in A(x),X=x\}=1-\theta_A(x).
$$
Thus, the path through $\Pi(x)$ is determined by the realized response. The covariates determine the partition and the probability of each transition along that path, rather than assigning the observation deterministically to a branch as in an ordinary regression tree.

For a node label $A$, let $\mathcal P(A)$ denote its internal ancestor labels. This set depends only on the binary topology. The conditional probability assigned to $A(x)$ is
\begin{equation}
    P\{T\in A(x)\mid X=x\}
    =
    \prod_{B\in\mathcal P(A)}
    \theta_B(x)^{\mathbbm{1}\{A(x)\subseteq B^\ell(x)\}}
    \{1-\theta_B(x)\}^{\mathbbm{1}\{A(x)\subseteq B^r(x)\}} .
    \label{eq:deterministic-response-tree-mass}
\end{equation}
Write $p_A(x)=P\{T\in A(x)\mid X=x\}$. The child masses satisfy
\[
    p_{A^\ell}(x)=p_A(x)\theta_A(x),\qquad
    p_{A^r}(x)=p_A(x)\{1-\theta_A(x)\},\qquad
    p_{A^\ell}(x)+p_{A^r}(x)=p_A(x).
\]
Since $p_{A_{0,1}}(x)=1$, summing over each level gives
\[
    \sum_{B\in\mathcal N_{m+1}}p_B(x)
    =\sum_{A\in\mathcal N_m}\{p_{A^\ell}(x)+p_{A^r}(x)\}
    =\sum_{A\in\mathcal N_m}p_A(x)
    =\cdots=p_{A_{0,1}}(x)=1.
\]
Thus the nonnegative node masses are normalized and consistent under refinement of the tree partitions. Both the partition $\Pi(x)$ and the branching probabilities $\theta_A(x)$ are unknown functions to be estimated from the data.

\subsection{Survival Data}
Censored survival outcomes arise when subjects are inspected intermittently, leave a study before the event occurs, or reach the end of follow-up without an observed event. For subject $i$, let $T_i \in \mathcal{S}=\mathbb{R}^{+}$ denote the latent event time of interest, and let $X_i \in \mathbb{R}^{d}$ denote a $d$-dimensional vector of baseline or subject-specific covariates. The covariate vector $X_i$ is assumed to be fully observed. The censoring or inspection process is assumed to be independent of $T_i$ conditional on $X_i$.

The observed survival outcome is represented by its endpoints, written as $(L_i,R_i]$, where $0\leq L_i\leq R_i\leq\infty$. When $L_i<R_i$, the observation specifies $T_i\in(L_i,R_i]$; when $L_i=R_i$, it specifies the exact event time $T_i=L_i$. Right censoring at time $C_i$ gives $(L_i,R_i)=(C_i,\infty)$. For current-status data, a single inspection at $C_i$ gives $(L_i,R_i]=(0,C_i]$ if $T_i\leq C_i$ and $(C_i,\infty)$ otherwise. Under mixed censoring, different subjects may contribute different observation types.
For a sample of $n$ independent subjects, the observed data are given by
$$
    \{(y_i,x_i): i=1,\ldots,n\},
    \qquad
    y_i=(l_i,r_i),
$$
where $(l_i,r_i,x_i)$ is the realization of $(L_i,R_i,X_i)$. Lowercase arguments are used when the likelihood is evaluated at the observed data.

\subsection{Finite Regression Tree and Likelihood Modeling}
\label{section-FNRT}
For computation, the tree is truncated at a finite depth $M$. Both the subject-specific partition and the probability allocation are allowed to depend on covariates.

The response distribution is modeled on the collapsed space $\overline{\mathcal S}_\tau=[0,\tau]\mathbin{\dot\cup}\{\dagger\}$, where $\dagger$ represents $T>\tau$ and the body gives no mass to zero. The body geometry is generated by a depth-$M$ binary tree with root interval $A_{0,1}(x)=[0,\tau]$. Let $\mathcal N_m$ denote the fixed node labels at level $m$, corresponding to the intervals in $\mathcal N_m(x)$. The internal and terminal label sets are
\[
    \mathcal I_M=\bigcup_{m=0}^{M-1}\mathcal N_m,
    \qquad \mathcal V_M=\mathcal N_M,
    \qquad J=|\mathcal V_M|=2^M.
\]
For subject $i$, the interval associated with label $A$ is $A(x_i)$. Thus the labels and ancestral paths are fixed, while the interval endpoints depend on the covariates.

Let $\nu_R$ be the rightmost body terminal label. In the probability tree, the tail is attached to the rightmost node at every level. Each node on that path represents its body interval together with $\dagger$; all other nodes represent their body intervals alone. Below level $M$, the rightmost node is split at the fixed endpoint $\tau$ into $\nu_R(x)$ and $\dagger$. Denote this additional internal label by $A_\dagger$, and put
\[
    \overline{\mathcal I}_M=\mathcal I_M\cup\{A_\dagger\},
    \qquad \overline{\mathcal V}_M=\mathcal V_M\cup\{\dagger\}.
\]
The cut head has $J-1$ outputs indexed by $\mathcal I_M$, whereas the probability head has $J$ outputs indexed by $\overline{\mathcal I}_M$. The extra split has a fixed location and needs no cut proportion. Widths and cut proportions always refer to the finite body geometry. A branching probability on the rightmost path conditions on the corresponding body interval together with the tail.

For a body terminal label $\nu\in\mathcal V_M$, the original depth-$M$ path $\mathcal P(\nu)$ gives the mass before the additional split,
$$
    \widetilde\pi_\nu(x_i)
    =
    \prod_{B\in \mathcal{P}(\nu)}
    \theta_B(x_i)^{\mathbbm{1}\{\nu(x_i)\subseteq B^\ell(x_i)\}}
    \{1-\theta_B(x_i)\}^{\mathbbm{1}\{\nu(x_i)\subseteq B^r(x_i)\}}.
$$

For $\nu\neq\nu_R$, the mass is already a body terminal mass. The final split gives
\begin{equation}
\begin{aligned}
    \pi_\nu(x)&=\widetilde\pi_\nu(x), &&\nu\neq\nu_R,\\
    \pi_{\nu_R}(x)&=\widetilde\pi_{\nu_R}(x)\theta_{A_\dagger}(x),&\qquad
    s(x)&=\widetilde\pi_{\nu_R}(x)\{1-\theta_{A_\dagger}(x)\}.
\end{aligned}
\label{eq:body-tail-masses}
\end{equation}
The normalization before the final split gives $\sum_{\nu\in\mathcal V_M}\widetilde\pi_\nu(x)=1$. Hence
\begin{align*}
    \sum_{\nu\in\mathcal V_M}\pi_\nu(x)+s(x)
    &=\sum_{\nu\ne\nu_R}\widetilde\pi_\nu(x)
      +\widetilde\pi_{\nu_R}(x)
         [\theta_{A_\dagger}(x)+\{1-\theta_{A_\dagger}(x)\}]\\
    &=\sum_{\nu\in\mathcal V_M}\widetilde\pi_\nu(x)=1,
\end{align*}
where $s(x)=P(T>\tau\mid X=x)$. Equivalently, all masses in $\overline{\mathcal V}_M$ are products along paths of the expanded probability tree; write $\overline{\mathcal P}(\nu)$ for these paths. Their lengths are at most $M+1$.

A finite tree specifies the probability mass assigned to each terminal body interval and to the tail. A uniform distribution is assumed within each terminal body interval; the distribution within the tail is left unspecified. For a censored observation $y_i=(l_i,r_i)$ with $l_i<r_i$, the contribution from terminal node $\nu$ is the fraction of $\nu(x_i)$ overlapping $(l_i,r_i]$:
$$
    q_\nu(l_i,r_i;x_i)
    =
    \frac{
        \left[\min\{b_\nu(x_i),r_i\}-\max\{a_\nu(x_i),l_i\}\right]_+
    }{
        W_\nu(x_i)
    },
$$
where $[u]_+=\max(u,0)$, and right censoring is represented by $r_i=\infty$. For an exact event time, $l_i=r_i=t_i$, and the contribution is the uniform density on the terminal interval containing $t_i$:
$$
    q_\nu(t_i,t_i;x_i)
    =
    \frac{\mathbbm{1}\{t_i\in\nu(x_i)\}}{W_\nu(x_i)}.
$$

The likelihood contribution for subject $i$ is
$$
    p(y_i\mid x_i)
    =
    \sum_{\nu\in\mathcal V_M}
    \pi_\nu(x_i)
    q_\nu(l_i,r_i;x_i)+s(x_i)\mathbbm{1}\{r_i=\infty\}.
$$

This likelihood applies to exact, right-censored, current-status, and interval-censored observations whose finite endpoints are at most $\tau$. Every observation set must contain either the entire tail or none of it. Information beyond $\tau$ is therefore recorded only as $T>\tau$; exact times beyond $\tau$ are represented by $(\tau,\infty)$. No density or partial-interval probability within the tail is specified. Direct maximization over all branch-specific tree parameters is difficult. The covariate-dependent parameters are therefore represented by a multilayer perceptron (MLP) with two output heads, allowing the observed-data likelihood maximum to be approximated through gradient descent. Estimation is based on the penalized negative log-likelihood introduced below.

\subsection{Neural Parameterization and Partition Construction}
\label{section-neural-parameterization}
In the finite regression tree of Section \ref{section-FNRT}, both the response partition and the probability allocation depend on covariates. These quantities are parameterized by a two-head MLP that maps the covariates to the cut proportion $\rho_A(x_i,w)$ and the branching probability $\theta_A(x_i,w)$. The body topology is fixed at a chosen depth $M$, with the additional probability split at $\tau$. Split locations are estimated from the observed-data likelihood and may vary with covariates, allowing both terminal interval widths and terminal probabilities to differ across subjects.

Let $w$ collect the weights and biases of the neural network. For subject $i$, the network first maps the covariates to a shared hidden representation,
$$
    \bm h_i(w)=g_w(x_i).
$$
When no shared hidden layers are used, $g_w$ is the identity map. Two separate MLP branches produce cut logits $\eta_A^{(\rho)}(x_i,w)$ for $A\in\mathcal I_M$ and branch logits $\eta_A^{(\theta)}(x_i,w)$ for $A\in\overline{\mathcal I}_M$. Write $\operatorname{expit}(z)=(1+e^{-z})^{-1}$ for the logistic function. The first head specifies the covariate-dependent cut proportion,
\begin{equation}
    \label{eq:cut_prop}
    \rho_A(x_i,w)
    =
    \epsilon_\rho + (1-2\epsilon_\rho) \operatorname{expit}\{\eta_A^{(\rho)}(x_i,w)\},
    \qquad 0 < \epsilon_\rho < 1/2,
\end{equation}
so that $\rho_A(x_i,w)\in[\epsilon_\rho,1-\epsilon_\rho]$. This proportion is used to allocate the available width after space has been reserved for all descendant terminal intervals.

Fix a reservation fraction $0\leq\beta<1$ and define
\begin{equation}
    h_M=\frac{\beta\tau}{2^M},\qquad
    r_{m,M}=2^{M-m-1}h_M,\qquad m=0,\ldots,M-1.
    \label{eq:width-reservation}
\end{equation}
Here $h_M$ is the minimum reserved width of a terminal interval, and $r_{m,M}$ is the width reserved for each child of a level-$m$ node, since that child has $2^{M-m-1}$ terminal descendants. The constant $\beta$ specifies the total reserved fraction of the root interval and is fixed before fitting. 

Including the network argument $w$ in the endpoint and width functions, the split point of a level-$m$ node is
\begin{equation}
    c_A(x_i,w)=a_A(x_i,w)+r_{m,M}
      +\rho_A(x_i,w)\{W_A(x_i,w)-2r_{m,M}\}.
    \label{eq:general-tree-cut}
\end{equation}
For $A\in\mathcal N_m$, define its unreserved width by
\[
    E_A(x_i,w)=W_A(x_i,w)-2^{M-m}h_M.
\]
The split rule gives the exact recursions
\begin{align*}
    W_{A^\ell}(x_i,w)&=r_{m,M}+\rho_A(x_i,w)E_A(x_i,w),\\
    W_{A^r}(x_i,w)&=r_{m,M}+\{1-\rho_A(x_i,w)\}E_A(x_i,w),\\
    E_{A^\ell}(x_i,w)&=\rho_A(x_i,w)E_A(x_i,w),\\
    E_{A^r}(x_i,w)&=\{1-\rho_A(x_i,w)\}E_A(x_i,w).
\end{align*}
At the root, $E_{A_{0,1}}=(1-\beta)\tau>0$. Along a path $A^{[0]}=A_{0,1},\ldots,A^{[m]}=A$, let $q_j^{(\rho)}=\rho_{A^{[j]}}(x_i,w)$ for a left transition and $q_j^{(\rho)}=1-\rho_{A^{[j]}}(x_i,w)$ for a right transition. Then
\[
    E_A(x_i,w)=(1-\beta)\tau\prod_{j=0}^{m-1}q_j^{(\rho)}
       \geq(1-\beta)\tau\epsilon_\rho^m>0.
\]
Consequently, recursion from $[0,\tau]$ defines the terminal partition and gives
\begin{equation}
    W_A(x_i,w)>2^{M-m}h_M
    \qquad\text{for every node at level }m=0,\ldots,M.
    \label{eq:reservation-invariant}
\end{equation}
In particular, every terminal width exceeds $h_M$. To bound the effective width ratios, suppress $(x_i,w)$ and put $u_A=2r_{m,M}/W_A\in[0,1)$. Then
\[
    \frac{W_{A^\ell}}{W_A}=\frac{u_A}{2}+(1-u_A)\rho_A,
    \qquad
    \frac{W_{A^r}}{W_A}=\frac{u_A}{2}+(1-u_A)(1-\rho_A).
\]
For either child $B$, these convex combinations satisfy
\[
    \epsilon_\rho
    =u_A\epsilon_\rho+(1-u_A)\epsilon_\rho
    \leq\frac{W_B}{W_A}
    \leq u_A(1-\epsilon_\rho)+(1-u_A)(1-\epsilon_\rho)
    =1-\epsilon_\rho.
\]
Thus $W_\nu=\tau\prod_{j=0}^{M-1}(W_{A^{[j+1]}}/W_{A^{[j]}})\geq\tau\epsilon_\rho^M>0$ for each terminal node. Setting $\beta=0$ gives $r_{m,M}=0$ and $c_A=a_A+\rho_AW_A$, recovering the proportional split in Section~\ref{section-regression-tree}. For $\rho_A=1/2$,
\[
    c_A=a_A+r_{m,M}+\tfrac12(W_A-2r_{m,M})
       =a_A+\tfrac12W_A,\qquad
    W_{A^\ell}=W_{A^r}=\tfrac12W_A.
\]

The second head determines the left-branching probability for every $A\in\overline{\mathcal I}_M$, including the final body--tail split,
\begin{equation}
    \label{eq:branch_prob}
    \theta_A(x_i,w)
    =
    \operatorname{expit}\{\eta_A^{(\theta)}(x_i,w)\},
    \qquad 0<\theta_A(x_i,w)<1.
\end{equation}

\begin{figure}[ht]
\centering
\begin{adjustbox}{max width=0.96\textwidth}
\begin{tikzpicture}[
  >=Latex,
  font=\sffamily\small,
  box/.style={
    draw=black!70,
    rounded corners=2pt,
    thick,
    align=center,
    inner xsep=5pt,
    inner ysep=4pt,
    fill=white
  },
  input/.style={box, minimum width=38mm, minimum height=9mm},
  hidden/.style={box, minimum width=46mm, minimum height=10mm},
  mlp/.style={box, minimum width=42mm, minimum height=12mm},
  output/.style={box, minimum width=46mm, minimum height=14mm},
  final/.style={box, minimum width=100mm, minimum height=15mm},
  group/.style={draw=black!25, rounded corners=3pt, fill=black!5, inner sep=7pt},
  label/.style={font=\sffamily\footnotesize\bfseries, text=black!65, fill=black!5, inner xsep=2pt, inner ysep=1pt},
  arr/.style={->, thick, draw=black!70},
  split/.style={->, thick, draw=black!70, rounded corners=4pt}
]

\node[input] (x) {\textbf{Covariates}\\[-1mm] $x_i$};
\node[hidden, right=18mm of x] (h) {\textbf{Shared Embedding MLP}\\[-1mm] $\bm h_i=g_w(x_i)$};

\node[mlp, below=22mm of h, xshift=-34mm] (rhomlp) {\textbf{Partition MLP}\\[-.5mm]
$\bm h_i\mapsto\eta_A^{(\rho)}$};
\node[output, below=9mm of rhomlp] (rhoout) {$\eta_A^{(\rho)}\mapsto \rho_A$, $A\in\mathcal I_M$\\[-.5mm]
$\rho_A=\epsilon_\rho+(1-2\epsilon_\rho)\operatorname{expit}(\eta_A^{(\rho)})$};

\node[mlp, below=22mm of h, xshift=34mm] (thetalp) {\textbf{Branching-Probability MLP}\\[-.5mm]
$\bm h_i\mapsto\eta_A^{(\theta)}$};
\node[output, below=9mm of thetalp] (thetaout) {$\eta_A^{(\theta)}\mapsto \theta_A$, $A\in\overline{\mathcal I}_M$\\[-.5mm]
$\theta_A=\operatorname{expit}(\eta_A^{(\theta)})$};

\node[final, below=24mm of rhoout |- thetaout.south, xshift=34mm] (tree) {\textbf{Regression Tree Node}\\[1mm]
$\rho_A$ defines $(A^\ell,A^r)$, \qquad
$\theta_A$ allocates mass, including the tail split};

\draw[arr] (x) -- (h);
\draw[split] ([xshift=-10mm]h.south) -- ++(0,-8mm) -| (rhomlp.north);
\draw[split] ([xshift=10mm]h.south) -- ++(0,-8mm) -| (thetalp.north);
\draw[arr] (rhomlp) -- (rhoout);
\draw[arr] (thetalp) -- (thetaout);
\draw[split] (rhoout.south) -- ++(0,-8mm) -| ([xshift=-25mm]tree.north);
\draw[split] (thetaout.south) -- ++(0,-8mm) -| ([xshift=25mm]tree.north);

\begin{scope}[on background layer]
  \node[group, fit=(x)(h)] (genc) {};
  \node[group, fit=(rhomlp)(rhoout)] (gpart) {};
  \node[group, fit=(thetalp)(thetaout)] (gtheta) {};
  \node[group, fit=(tree)] (gtree) {};
\end{scope}
\node[label, anchor=south] at ([yshift=.7mm]genc.north) {Shared Encoder};
\node[label, anchor=south] at ([yshift=.7mm]gpart.north) {Covariate-Dependent Partition};
\node[label, anchor=south] at ([yshift=.7mm]gtheta.north) {Branching Probabilities};
\node[label, anchor=south] at ([yshift=.7mm]gtree.north) {Neural Regression Tree Distribution};

\end{tikzpicture}
\end{adjustbox}
\caption{Neural Regression Tree construction. The cut and probability heads have $J-1$ and $J$ outputs, respectively; the additional probability split separates the last body interval from the tail at $\tau$.}
\label{fig:nn-structure}
\end{figure}

Given the deterministic branching probabilities $\{\theta_A(x_i,w):A\in\overline{\mathcal I}_M\}$, the mass before the additional split is obtained along the original depth-$M$ path,
$$
    \widetilde\pi_{\nu}(x_i,w)
    =
    \prod_{A\in \mathcal{P}(\nu)}
    \theta_A(x_i,w)^{\mathbbm{1}\{\nu(x_i,w)\subseteq A^\ell(x_i,w)\}}
    \{1-\theta_A(x_i,w)\}^{\mathbbm{1}\{\nu(x_i,w)\subseteq A^r(x_i,w)\}} .
$$
The body and tail masses are then obtained from \eqref{eq:body-tail-masses}. Let $q_\nu(l_i,r_i;x_i,w)$ be the contribution defined in Section \ref{section-FNRT}, evaluated on the covariate-dependent terminal intervals $\nu(x_i,w)$. The finite-tree likelihood is then
\begin{equation}
    p(y_i\mid x_i,w)
    =
    \sum_{\nu\in\mathcal V_M}
    \pi_{\nu}(x_i,w)
    q_{\nu}(l_i,r_i;x_i,w)+s(x_i,w)\mathbbm{1}\{r_i=\infty\}.
\label{eq:ll_contribution}
\end{equation}
\subsection{Penalized Likelihood Estimation}
\label{section-nrt-estimation}

The partition and branching-probability heads are estimated jointly by minimizing a penalized negative log-likelihood. The likelihood contribution in \eqref{eq:ll_contribution} gives
\begin{equation}
    \mc L_{\mathrm{NLL},n}(w)
    =-\frac{1}{n}\sum_{i=1}^n\log p(y_i\mid x_i,w).
    \label{eq:nll}
\end{equation}
This criterion uses a terminal density for an exact observation and an interval probability for a censored observation. Thus, exact, right-censored, interval-censored, and current-status observations inform the same conditional distribution through a common likelihood.

Estimating interval widths allows the response resolution to vary with covariates. However, likelihood maximization may favor narrow intervals around exact observations. For an exact event time $t_i$ in terminal interval $\nu_i(x_i,w)$, the negative log-likelihood contribution is
\begin{equation}
    -\log p(y_i\mid x_i,w)
    =-\log\pi_{\nu_i}(x_i,w)+\log W_{\nu_i}(x_i,w).
    \label{eq:likelihood-width-effect}
\end{equation}
For a fixed terminal mass, decreasing the interval width increases the likelihood as long as the observed event remains inside the interval. The positive cut margin and minimum terminal width limit this contraction. A geometric penalty is introduced to control departures from equal terminal widths.

Define the normalized terminal widths by
\[
    d_\nu(x,w)=\frac{W_\nu(x,w)}{\tau},
    \qquad \sum_{\nu\in\mathcal V_M}d_\nu(x,w)=1.
\]
The subject-specific partition penalty is
\begin{equation}
    \mathcal R_M(x,w)
    =-\frac{2}{MJ}\sum_{\nu\in\mathcal V_M}
       \log\{Jd_\nu(x,w)\}
    =\frac{2}{M}\operatorname{KL}\{\mathbf u_J\,\|\,\mathbf d(x,w)\},
    \label{eq:partition-penalty}
\end{equation}
where $\mathbf u_J=(1/J,\ldots,1/J)$ and $\mathbf d(x,w)$ is the vector of normalized terminal body widths. The tail has no width and is excluded from this penalty. This penalty is nonnegative, vanishes for equal terminal widths, and increases as any normalized width approaches zero. Its dependence on widths alone leaves the branching probabilities free to allocate unequal terminal masses. The fitted distribution can therefore remain nonuniform even when the penalty favors a nearly equal-width partition.

The resulting estimation criterion is
\begin{equation}
    \widehat w
    \in\argmin_w\mathcal J_n(w),
    \qquad
    \mathcal J_n(w)
    =\mc L_{\mathrm{NLL},n}(w)
     +\lambda_n\frac{1}{n}\sum_{i=1}^n\mathcal R_M(x_i,w),
    \qquad \lambda_n=\frac{\alpha}{\sqrt n},
    \label{eq:objective}
\end{equation}
where $\alpha\geq0$ is a tuning parameter and $n$ is the number of subjects used in the fit. The parameter $\alpha$ controls the tradeoff between likelihood fit and departures from equal terminal widths. Setting $\alpha=0$ gives unpenalized adaptive likelihood estimation. Larger values of $\alpha$ favor more even widths while allowing departures supported by the likelihood. For a fixed $\alpha$, the penalty coefficient decreases as the fitting sample grows. Both the probability allocation and the covariate-dependent partition are estimated through this penalized likelihood.

For a minibatch $\mc B$, the stochastic objective is
\begin{equation}
    \begin{aligned}
    \mathcal J_{\mc B}(w)
    &=\mc L_{\mathrm{NLL},\mc B}(w)
      +\frac{\lambda_n}{|\mc B|}
         \sum_{i\in\mc B}\mathcal R_M(x_i,w),\\
    \mc L_{\mathrm{NLL},\mc B}(w)
    &=-\frac{1}{|\mc B|}\sum_{i\in\mc B}\log p(y_i\mid x_i,w).
    \end{aligned}
    \label{eq:minibatch-nll}
\end{equation}
The objective is optimized using minibatch Adam. Write $w=(w_s,w_\theta,w_\rho)$, where $w_s$ contains any shared hidden-layer parameters and $w_\theta$ and $w_\rho$ contain the parameters specific to the branching probabilities and cut proportions, respectively. After a brief initialization phase with balanced cuts, all active parameters are updated using gradients of the same penalized loss. The probability allocation and partition are thus estimated jointly. Unpenalized held-out NLL is used for early stopping and cross-validation, so candidate models are evaluated on observations not used in fitting. The penalty is applied only during estimation. Tree depth, network architecture, and penalty strength are selected by cross-validation. The selected model is then refitted on the complete training sample for a duration determined from the cross-validation fits. Algorithm~\ref{alg:training} summarizes the procedure.

\begin{algorithm}[H]
\caption{Penalized likelihood training for the adaptive Neural Regression Tree}
\label{alg:training}
\begin{algorithmic}[1]
\Require Fitting and validation data, depth $M$, horizon $\tau$, network architecture, cut constraints, penalty strength $\alpha$, and stopping rule
\State Set $\lambda_n=\alpha/\sqrt n$, using the size $n$ of the fitting sample.
\State Initialize $w=(w_s,w_\theta,w_\rho)$ with balanced cuts and branching probabilities.
\State Perform a brief initialization phase with $\rho_A=1/2$ held fixed.
\For{each subsequent epoch}
    \For{each shuffled minibatch $\mc B$}
        \State Compute the $J-1$ cut proportions, $J$ branching probabilities, body intervals, and body and tail masses.
        \State Evaluate the observed-data NLL and width penalty in \eqref{eq:minibatch-nll}.
        \State Jointly update $(w_s,w_\theta,w_\rho)$ by an Adam step on $\mathcal J_{\mc B}(w)$.
    \EndFor
    \State Evaluate unpenalized validation NLL and retain the best weights.
    \State Stop when the validation NLL satisfies the stopping rule.
\EndFor
\State Restore and return the best weights.
\end{algorithmic}
\end{algorithm}

For regression, split locations and branching probabilities are estimated jointly as functions of the covariates under the penalized observed-data likelihood. Within the chosen binary topology, a single adaptive tree is used to represent the conditional distribution. This avoids the additional computation required to fit multiple trees and average their predictions in a forest.

Fixing $\rho_A(x,w)=1/2$ at every body internal node gives the fixed-partition NRT, for which all terminal widths equal $\tau/J$ and $\mathcal R_M(x,w)=0$. In the simulation study, the original finite-support versions of the fixed-partition and adaptive NRT are compared, with each selected by held-out NLL. A further fixed-partition control is fitted using the tree depth and network architecture selected for the adaptive model. Its fitting duration is selected separately by held-out NLL. This comparison assesses the contribution of the learned partition at the same tree depth and network architecture.


\section{Theoretical Properties}
\label{sec:theory}

The approximation and estimation properties of the NRT are studied under general covariate-dependent cuts. Under the assumptions stated below, the analysis establishes validity of the induced conditional distribution, uniform approximation of conditional distribution functions, identification under exact and censored observation schemes, and consistency of the penalized likelihood estimator over increasing model classes. Technical lemmas and all proofs are given in Section~S1 of the Supplementary Material. Throughout this section, the horizon $0<\tau<\infty$ is fixed.

The target is the conditional law of $Y=T$ for $T\leq\tau$ and $Y=\dagger$ otherwise, denoted by $\bar P_*(\cdot\mid x)$. Let $F_*(t\mid x)=P(T\leq t\mid X=x)$ for $t\in[0,\tau]$, and define the true conditional tail probability by
\[
    s_*(x):=P(T>\tau\mid X=x)=1-F_*(\tau\mid x).
\]
The tail probability may vary with $x$ and may equal zero or one.

\begin{assum}
\label{assum:nrt_target}
The covariate space $\mathcal X\subset\mathbb R^d$ is compact. The conditional distribution function $F_*(t\mid x)$ is jointly continuous on $[0,\tau]\times\mathcal X$ and satisfies $F_*(0\mid x)=0, x\in\mathcal X$.
\end{assum}

\subsection{Tree Geometry and the Partition Penalty}

The body depth $M\geq1$, body terminal count $J$, node labels, endpoints, and widths are as defined in Section~\ref{sec:NRT}. The horizon $\tau$, cut margin $\epsilon_\rho$, and reservation fraction $\beta$ are common to all trees under comparison.

For the approximation analysis, the cut logits $\eta_A^{(\rho)}$, $A\in\mathcal I_M$, and branch logits $\eta_A^{(\theta)}$, $A\in\overline{\mathcal I}_M$, are allowed to be arbitrary elements of $C(\mathcal X)$, the space of continuous real-valued functions on $\mathcal X$. The cut proportions $\rho_A$ and branch probabilities $\theta_A$ are obtained through the transformations in \eqref{eq:cut_prop} and \eqref{eq:branch_prob}. Functions written without $w$ refer to these general continuous tree parameters. The weight argument $w$ is retained when the logits are produced by the neural network of Section~\ref{section-neural-parameterization}. The same convention applies to node intervals, terminal masses, and the partition penalty.

The width reservation and split rule are given by \eqref{eq:width-reservation} and \eqref{eq:general-tree-cut}, with $w$ omitted under this convention. An admissible cut system consists of continuous cut logits transformed by \eqref{eq:cut_prop} and used in this split rule. Both $\beta$ and $\epsilon_\rho$ are held fixed as $M$ increases; the results include the proportional-split case $\beta=0$.

The terminal path-product masses $\pi_\nu(x)$ and uniform within-node distributions from Section~\ref{section-FNRT} define the body density and distribution function
\begin{align}
    f_M(t\mid x)
    &=\sum_{\nu\in\mathcal V_M}
      \frac{\pi_\nu(x)}{W_\nu(x)}
      \mathbbm{1}\{t\in\nu(x)\}, \label{eq:finite-tree-density}\\
    F_M(t\mid x)
    &=\sum_{\nu\in\mathcal V_M}\pi_\nu(x)G_\nu(t,x),
      \qquad
      G_\nu(t,x)=
      \frac{[\min\{t,b_\nu(x)\}-a_\nu(x)]_+}{W_\nu(x)}.
      \label{eq:finite-tree-cdf}
\end{align}
Densities are understood up to their values at interval endpoints. The conditional law on $\overline{\mathcal S}_\tau$ is $\bar P_M(dt\mid x)=f_M(t\mid x)\,dt+s_M(x)\delta_\dagger(dt)$, where $s_M(x)=1-F_M(\tau\mid x)$. The function $F_M$ is zero for $t<0$ and is specified on $[0,\tau]$; no CDF values at finite times beyond $\tau$ are assigned. The convention $F_M(\infty\mid x)=1$ represents the whole collapsed space.

For an internal node $A$ at level $m$, its effective left-child width proportion is
\begin{equation}
    \widetilde\rho_A(x)
      =\frac{W_{A^\ell}(x)}{W_A(x)}
      =\frac{r_{m,M}}{W_A(x)}
       +\left\{1-\frac{2r_{m,M}}{W_A(x)}\right\}\rho_A(x).
    \label{eq:effective-cut-proportion}
\end{equation}

\begin{lemma}
\label{lem:nrt_geometry_penalty}
For every admissible cut system and every collection of branch probabilities in $(0,1)$, the split rule in \eqref{eq:general-tree-cut} and the terminal body mixture in \eqref{eq:finite-tree-density}, together with the tail mass, define a conditional probability distribution on $\overline{\mathcal S}_\tau$ with body CDF \eqref{eq:finite-tree-cdf}. Its terminal widths satisfy
\begin{equation}
    \underline\Delta_M
    :=\max\{h_M,\tau\epsilon_\rho^M\}
    \leq W_\nu(x)\leq
    \overline\Delta_M:=\tau(1-\epsilon_\rho)^M.
    \label{eq:general-tree-mesh}
\end{equation}
In particular, $f_M(t\mid x)\leq\underline\Delta_M^{-1}$ for almost every $t$. The partition penalty in \eqref{eq:partition-penalty} satisfies
\begin{equation}
    0\leq\mathcal R_M(x)
    \leq C_\rho:=2\log\{(2\epsilon_\rho)^{-1}\},
    \label{eq:partition-penalty-bound}
\end{equation}
uniformly in $M$, $x$, and the tree parameters. If $\beta>0$, it also satisfies $\mathcal R_M(x)\leq 2M^{-1}\log(\beta^{-1})$.
\end{lemma}

The lemma provides a finite upper bound on exact-observation densities at each fixed depth and a uniform bound on the geometric penalty across depths. The penalty depends on interval widths; the terminal masses remain free to describe nonuniform conditional distributions.

The next result expresses the penalty through the level sets $\mathcal N_m$, $m<M$, defined in Section~\ref{section-FNRT}. 

\begin{proposition}
\label{prop:nrt_level_penalty}
For the effective proportions in \eqref{eq:effective-cut-proportion},
\begin{equation}
    \mathcal R_M(x)
    =\frac2M\sum_{m=0}^{M-1}2^{-m}\sum_{A\in\mathcal N_m}
        \operatorname{KL}\{\operatorname{Bern}(1/2)
                     \|\operatorname{Bern}(\widetilde\rho_A(x))\}.
    \label{eq:level-penalty-decomposition}
\end{equation}
Consequently,
\begin{align}
    \mathcal R_M(x)&\geq\frac4M\sum_{m=0}^{M-1}2^{-m}
          \sum_{A\in\mathcal N_m}(\widetilde\rho_A(x)-1/2)^2,
          \label{eq:local-width-control}\\
    \|\mathbf d(x)-\mathbf u_J\|_1^2&\leq M\mathcal R_M(x).
          \label{eq:terminal-width-control}
\end{align}
The penalty is zero if and only if all effective splits are balanced.
\end{proposition}

\subsection{Approximation under General Covariate-Dependent Cuts}

Let $\mathcal T_M$ be the class of conditional distribution functions in \eqref{eq:finite-tree-cdf} with the continuous cut functions above and continuous branch probabilities in $(0,1)$. Write $\|\cdot\|_\infty$ for the uniform norm over $\mathcal X$ for functions of $x$, and over $[0,\tau]\times\mathcal X$ for functions of $(t,x)$. Let $P_X$ denote the marginal distribution of $X$.

\begin{myTheo}
\label{thm:rt_full_support}
Suppose Assumption~\ref{assum:nrt_target} holds. For any sequence of admissible continuous cut systems, one at each depth $M$, and every $\varepsilon>0$, there exist a depth $M$ and continuous branch probabilities such that the resulting $F_M\in\mathcal T_M$ satisfies
\begin{equation}
    \|F_M-F_*\|_\infty<\varepsilon.
    \label{eq:rt_uniform_approximation}
\end{equation}
\end{myTheo}

This approximation property is next extended to the two-head neural parameterization. Neural-network sieve approximation is discussed by \citet{shen_asymptotic_2023}. Here both the cut functions and the branching probabilities must be approximated.

\begin{assum}
\label{assum:nrt_network_richness}
For each fixed $M$, the classes of joint pre-sigmoid outputs of the two heads are nested in a network complexity index $K$. All hidden-layer activations are continuous, and the hidden-layer activation admits uniform approximation of continuous scalar functions on $\mathcal X$. As $K$ increases, the architecture permits any finite collection of such scalar approximations to be represented simultaneously by the two heads, allowing the widths of the shared and private layers to increase. The cut and probability outputs are indexed by $\mathcal I_M$ and $\overline{\mathcal I}_M$, respectively, and each head has its own final affine output layer. Setting the weights and biases of the cut head's final layer to zero leaves the shared representation and the branch head free to approximate any finite collection of continuous branch logits uniformly as $K$ increases.
\end{assum}

\begin{myTheo}
\label{thm:nrt_global_approximation}
Under Assumptions~\ref{assum:nrt_target} and \ref{assum:nrt_network_richness}, for every $\varepsilon>0$ there exist a depth $M$, a network complexity $K$, and weights $w$ such that
\begin{equation}
    \|F_M(\cdot\mid\cdot,w)-F_*\|_\infty<\varepsilon.
    \label{eq:nrt-uniform-cdf-approximation}
\end{equation}
\end{myTheo}

\subsection{Identification from Exact and Censored Observations}
\label{subsec:identification-consistency}

Let $\xi=(M,K,w)$ index an NRT representation. Use $F_\xi(t\mid x)=F_M(t\mid x,w)$, $f_\xi(t\mid x)=f_M(t\mid x,w)$, and $\mathcal R_\xi(x)=\mathcal R_M(x,w)$ for its body CDF, body density, and partition penalty. Write $s_\xi(x)=1-F_\xi(\tau\mid x)$ and $\bar P_\xi(dt\mid x)=f_\xi(t\mid x)\,dt+s_\xi(x)\delta_\dagger(dt)$ for the tail mass and collapsed conditional law. For one subject, write $O=(X,D,L,R)$, where $D$ records the assigned observation design, the realized observation type, and, for an inspection design, the complete inspection schedule. For a candidate conditional CDF $F$ with density $f$, the response-dependent likelihood factor from Section~\ref{section-FNRT} is
\begin{equation}
    Q_F(O)=
    \begin{cases}
        f(L\mid X), & L=R,\\
        F(R\mid X)-F(L\mid X), & L<R,
    \end{cases}
    \qquad F(\infty\mid x)=1.
    \label{eq:observed-likelihood-factor}
\end{equation}
For an NRT representation, $Q_{F_\xi}(O_i)=p(y_i\mid x_i,w)$ in \eqref{eq:ll_contribution}. Let $P_F^O$ denote the observed-data law under $F$, and write $P_*^O$ for the true law.

\begin{assum}
\label{assum:observation_scheme}
The observations $O_1,\ldots,O_n$ are independent and identically distributed. All finite event, censoring, and inspection endpoints used in the likelihood lie in $[0,\tau]$. The observation mechanism depends on $T>\tau$ only through $Y=\dagger$, so a tail event can be recorded as $(\tau,\infty)$ but its exact time is not retained. Conditional on $X$, assignment to an observation design is independent of $T$. Within each assigned design, the censoring or inspection process is independent of $T$ conditional on $X$. For inspection designs, the full schedule, including its almost surely finite number of inspections, is recorded. The covariate distribution and observation mechanism are held fixed when comparing candidate event-time distributions. With respect to a common dominating measure, there is a function $h$ independent of $F$ such that
\[
    p_F^O(o)=h(o)Q_F(o).
\]
\end{assum}

\begin{assum}
\label{assum:identifying_support}
For $P_X$-almost every $x$, at least one recorded design subgroup has positive conditional probability and satisfies one of the following conditions, with all probabilities interpreted conditional on $X=x$ and subgroup membership:
\begin{enumerate}
    \item $Y$ is observed for every subject in the subgroup: event times up to $\tau$ are exact and $T>\tau$ is recorded otherwise;
    \item observations are right-censored, and $P(C\geq t\mid X=x)>0$ for every $t<\tau$;
    \item observations are obtained by current-status or interval-censoring inspections, and the union of the conditional supports of the recorded inspection times is dense in $[0,\tau]$.
\end{enumerate}
For a single observation design, the subgroup is the entire population.
\end{assum}

\begin{myTheo}
\label{thm:nrt_identifiability}
Under Assumptions~\ref{assum:observation_scheme} and \ref{assum:identifying_support}, if two NRT representations $\xi_1$ and $\xi_2$ satisfy $P_{F_{\xi_1}}^O=P_{F_{\xi_2}}^O$, then
\[
    F_{\xi_1}(t\mid x)=F_{\xi_2}(t\mid x)
    \quad\text{for every }t\in[0,\tau]
    \text{ and for }P_X\text{-almost every }x.
\]
\end{myTheo}

The theorem identifies the conditional distribution. Distinct cut functions, branch probabilities, and neural-network weights may represent the same distribution.

\subsection{Consistency of Penalized Likelihood Estimation}

Exact observations contribute a density to the likelihood. For consistency, CDF approximation is supplemented by approximation in integrated conditional KL divergence.

\begin{assum}
\label{assum:regular_target_density}
Fix $s_0\in(0,1)$ and let $\bar P_0=(1-s_0)U(0,\tau]+s_0\delta_\dagger$. For $P_X$-almost every $x$, the collapsed target law satisfies $\bar P_*(\cdot\mid x)\ll\bar P_0$, and
\begin{equation}
    A_0:=\int_{\mathcal X}\operatorname{KL}
       \{\bar P_*(\cdot\mid x)\|\bar P_0\}\,dP_X(x)<\infty.
    \label{eq:target-finite-kl}
\end{equation}
\end{assum}

This assumption requires an absolutely continuous body distribution but imposes neither continuity nor uniform upper or positive lower bounds on its density $f_*$. The tail mass $s_*(x)$ may be zero. The joint continuity of the body CDF in Assumption~\ref{assum:nrt_target} is retained. Write
\[
    a(\xi):=\int_{\mathcal X}\operatorname{KL}
       \{\bar P_*(\cdot\mid x)\|\bar P_\xi(\cdot\mid x)\}\,dP_X(x).
\]
The divergence is defined on the collapsed space and includes the tail term $s_*(x)\log\{s_*(x)/s_\xi(x)\}$, interpreted as zero when $s_*(x)=0$.

Balanced partitions provide KL approximations with zero partition penalty, as established in the Supplementary Material.

For estimation with sample size $n$, define a sieve of NRT representations by
\[
    \Xi_n=
    \bigcup_{\substack{1\leq M\leq M_n\\1\leq K\leq K_n}}
       \{(M,K,w):w\in\mathcal W_{M,K},\ \|w\|\leq C_{w,n}\},
\]
where $\mathcal W_{M,K}$ is the finite-dimensional Euclidean weight space of the specified architecture, and $M_n$, $K_n$, and $C_{w,n}$ are nondecreasing deterministic bounds. The complexity indices are integers, with one specified finite architecture for each pair $(M,K)$. The argument also allows finitely many architectures for each pair. The constants $\tau$, $\epsilon_\rho$, and $\beta$ are common to all sieve elements. The penalized criterion is indexed by $\xi$ because representations of the same CDF can have different partition penalties.

For an integrable function $g$ of the observed data, write $P_n g=n^{-1}\sum_{i=1}^n g(O_i)$ for the empirical average and $P_*g=\int g(o)\,dP_*^O(o)$ for expectation under the true law. Define
\[
    \ell_\xi(O)=\log Q_{F_\xi}(O),\qquad
    \ell_*(O)=\log Q_{F_*}(O),\qquad
    \overline{\mathcal R}_n(\xi)
       =\frac1n\sum_{i=1}^n\mathcal R_\xi(X_i),
\]
and define the relative observed-data log likelihood by
\begin{equation}
    \widetilde\ell_\xi(O)=\log\frac{Q_{F_\xi}(O)}{Q_{F_*}(O)}
            =\ell_\xi(O)-\ell_*(O).
    \label{eq:relative-observed-loglikelihood}
\end{equation}
These quantities are defined on the event $Q_{F_*}(O)>0$, which has probability one under the true observed-data law. The common mechanism factor in Assumption~\ref{assum:observation_scheme} cancels in the ratio. The analysis is based on this contrast, so separate integrability of $\ell_\xi$ and $\ell_*$ is not required. Distributional error is measured by
\[
    d_{\mathrm{CDF}}^2(F,G)=
       \int_{\mathcal X}\int_0^\tau
          \{F(t\mid x)-G(t\mid x)\}^2\,dt\,dP_X(x).
\]

Convergence in this metric gives $L^2$ consistency of the body CDF with respect to $dt\,dP_X(x)$. It does not by itself imply convergence of $F(\tau\mid x)$ or of the tail probability $s_F(x)=1-F(\tau\mid x)$. Additional endpoint support is used below to obtain tail-probability consistency. Write
\[
    \varkappa(F)=\operatorname{KL}(P_*^O\|P_F^O)
\]
for the observed-data KL divergence.

\Needspace{10\baselineskip}
\begin{assum}
\label{assum:sieve_consistency}
The sieves satisfy the following conditions:
\begin{enumerate}
    \item There is a deterministic sequence $\xi_n^\circ\in\Xi_n$ such that
    \begin{equation}
        \mathcal R_{\xi_n^\circ}(x)\equiv0,
        \qquad a(\xi_n^\circ)\longrightarrow0.
        \label{eq:slow-sieve-approximation}
    \end{equation}
    \item The uniform empirical error of the relative log likelihoods in \eqref{eq:relative-observed-loglikelihood} satisfies
    \begin{equation}
        e_n^{\mathrm{emp}}:=\sup_{\xi\in\Xi_n}|(P_n-P_*)\widetilde\ell_\xi|=o_p(1).
        \label{eq:sieve-ulln}
    \end{equation}
\end{enumerate}
\end{assum}

\begin{proposition}
\label{prop:nrt_slow_sieve}
Under Assumptions~\ref{assum:nrt_target}, \ref{assum:nrt_network_richness}, \ref{assum:observation_scheme}, and \ref{assum:regular_target_density}, there exist nondecreasing deterministic bounds $M_n,K_n,C_{w,n}\to\infty$ for which Assumption~\ref{assum:sieve_consistency} holds.
\end{proposition}

The construction in Proposition~\ref{prop:nrt_slow_sieve} may depend on the target and observation law. Thus, \eqref{eq:sieve-ulln} must still be verified for a specified architecture sequence. The technical approximation and likelihood-separation arguments are given in the Supplementary Material.

The empirical criterion corresponding to \eqref{eq:objective} is
\[
    \mathcal J_n(\xi)
       =-P_n\ell_\xi+\lambda_n\overline{\mathcal R}_n(\xi).
\]
Let $\widehat\xi_n\in\Xi_n$ be a measurable output satisfying the comparison condition with the sequence $\xi_n^\circ$ in Assumption~\ref{assum:sieve_consistency}:
\begin{equation}
    \mathcal J_n(\widehat\xi_n)
       \leq\mathcal J_n(\xi_n^\circ)+e_n^{\mathrm{opt}},
       \qquad e_n^{\mathrm{opt}}\geq0,\quad e_n^{\mathrm{opt}}=o_p(1).
    \label{eq:approximate-penalized-sieve}
\end{equation}
Approximate global minimization is sufficient for this condition, but only comparison with this one sequence is required. Convergence to a stationary point alone does not establish the condition.

\Needspace{15\baselineskip}
\begin{myTheo}
\label{thm:nrt_sieve_consistency}
Suppose Assumptions~\ref{assum:nrt_target}--\ref{assum:sieve_consistency} hold. Let $\lambda_n$ be any measurable, almost surely finite, nonnegative sequence, possibly depending on the data. If the comparison condition \eqref{eq:approximate-penalized-sieve} holds, then
\begin{equation}
    d_{\mathrm{CDF}}(F_{\widehat\xi_n},F_*)\xrightarrow{p}0.
    \label{eq:nrt-distributional-consistency}
\end{equation}
Moreover, $\varkappa(F_{\widehat\xi_n})+\lambda_n\overline{\mathcal R}_n(\widehat\xi_n)=o_p(1)$. If, for $P_X$-almost every $x$, a recorded subgroup of positive conditional probability observes $Y$ directly, has right censoring with $P(C\geq\tau\mid X=x)>0$, or has a recorded inspection at $\tau$ with positive conditional probability, then
\[
    \int_{\mathcal X}|s_{\widehat\xi_n}(x)-s_*(x)|^2\,dP_X(x)\xrightarrow{p}0.
\]
Subgroup conditioning is understood in these additional endpoint-support conditions.
\end{myTheo}

The estimator is allowed to range over the full adaptive NRT sieve, and balanced partitions provide the comparison sequence. This subclass has identically zero penalty and retains the required KL approximation property. Consistency therefore requires no decay condition on $\lambda_n$. In particular, data-dependent choices $\lambda_n=\alpha_n/\sqrt n$ are covered whenever the coefficient is finite and the stated optimization and sieve conditions hold.

\section{Simulation Studies}
\label{sec:simulation}

The numerical results and implementation details below are retained from the original finite-support NRT, which fixes the tail mass at zero. They have not been recomputed for the tail-node extension in Section~\ref{sec:NRT}.
The finite-sample performance of NRT was evaluated by estimating conditional survival functions under nonlinear covariate effects and different observation schemes. The comparisons assess the effect of learning the response partition, the performance of likelihood-based tuning, and the loss of accuracy as censoring increases. The simulation settings and results reported below follow the combinations displayed in Tables~\ref{tab:nrt_exact}--\ref{tab:nrt_current}.

\subsection{Data Generation}
Data are generated from a conditional survival model with six covariates. For $i=1,\ldots,n$, the covariate vector is $X_i=(X_{i1},\ldots,X_{i6})^\top$, where the six covariates are mutually independent and $X_{ij}\sim\mathrm{Unif}(-1,1)$. Only the first two covariates affect the event-time distribution. For a radius parameter $r>0$, define
\begin{equation}
    \eta_i=-\frac{X_{i1}^2+X_{i2}^2}{2r^2}.
    \label{eq:data-generation-predictor}
\end{equation}
The radius was set to $r=0.5$ for the proportional-odds model and $r=0.7$ for the other models. Thus, the covariate effects are nonlinear and radially symmetric in $(X_{i1},X_{i2})$, while $X_{i3},\ldots,X_{i6}$ are inactive. 

Five conditional event-time models are considered: proportional hazards (PH), lognormal, proportional odds (PO), and two multimodal cases. In the proportional-hazards scenario, let $\gamma_0=-\log(0.1)$. Event times are generated as
\begin{equation}
    T_i\mid X_i
    \sim \mathrm{Exponential}\left(\gamma_0\exp\{\eta_i\}\right),
    \label{eq:ph-data-generation}
\end{equation}
where the argument of the exponential distribution denotes its rate parameter. Equivalently,
\[
    \Pr(T_i>t\mid X_i)
    =\exp\{-\gamma_0 t\exp(\eta_i)\}.
\]
For lognormal event times,
\begin{equation}
    \log T_i\mid X_i
    \sim N(\eta_i,\sigma^2),
    \qquad \sigma=0.5.
    \label{eq:lognormal-data-generation}
\end{equation}
For the proportional-odds scenario, event times have finite support with endpoint $20$. Writing
\[
    s_i=\frac{\exp(-\eta_i/2)}{20},
\]
the conditional survival function is
\begin{equation}
    \Pr(T_i>t\mid X_i)
    =
    \begin{cases}
    \displaystyle
    \left[
    1+\left\{\frac{t}{(20-t)s_i}\right\}^{2}
    \right]^{-1},
    &0\leq t<20,\\[8pt]
    0,&t\geq 20.
    \end{cases}
    \label{eq:po-data-generation}
\end{equation}
Event times are sampled from this distribution by inverse transformation.

For the multimodal scenarios, two Weibull components with distinct shape parameters are mixed with equal weights. Let
\[
    (\kappa_1,\kappa_2)=(6,8),
    \qquad
    (\varpi_1,\varpi_2)=(0.5,0.5),
\]
and define the component-specific scales
\[
    s_{ik}
    =m_k\exp\left(-\frac{\eta_i}{\kappa_k}\right),
    \qquad k=1,2.
\]
Conditional on $X_i$, a latent component $Z_i$ is sampled with
$\Pr(Z_i=k\mid X_i)=\varpi_k$, followed by
\[
    T_i\mid X_i,Z_i=k
    \sim\mathrm{Weibull}(s_{ik},\kappa_k).
\]
Two sets of scale multipliers are considered,
\[
    (m_1,m_2)=(1,6)
    \qquad\text{and}\qquad
    (m_1,m_2)=(3,6).
\]
The first specification produces more widely separated modes, whereas the second produces closer modes. Marginalizing over $Z_i$ yields a mixture survival function,
\[
    \Pr(T_i>t\mid X_i)
    =
    \sum_{k=1}^2
    \varpi_k\exp\left\{-\left(\frac{t}{s_{ik}}\right)^{\kappa_k}\right\}.
\]

Censoring times are generated independently of $T_i$ and $X_i$. For each scenario and target right-censoring proportion $q\in\{0.2,0.5\}$, a common follow-up horizon $\tau_c$ is calibrated by deterministic quadrature over the covariate distribution. The covariate expectation is approximated on a $64\times64$ midpoint grid over $[-1,1]^2$. The horizon $\tau_c$ is then determined by bisection so that, for $C_i\sim\mathrm{Unif}(0,\tau_c)$,
\[
    E_{X,C}\{S_*(C\mid X)\}\approx q,
\]
where $S_*(t\mid x)=1-F_*(t\mid x)$ is the true survival function and the expectation is over both the covariates and the random censoring time. The corresponding probability is evaluated analytically for the proportional-hazards and lognormal models, by Simpson quadrature for the bounded proportional-odds model, and using the incomplete gamma function for the multimodal Weibull mixture. Censoring times are then sampled independently from the uniform distribution with the calibrated horizon.
Calibration is therefore performed at the population level, independently of the realized simulation sample. The calibrated horizon was saved during tuning and reused across replications. Here $q$ denotes $\Pr(R_i=\infty)$ under each censored observation scheme; it does not denote the proportion of all inexact observations under current-status or interval censoring. Realized right-censoring proportions fluctuate around $q$. For exact observations, $L_i=R_i=T_i$, and no censoring horizon is needed. For right-censored observations,
\[
    L_i=\min(T_i,C_i),
    \qquad
    R_i=
    \begin{cases}
    L_i,&T_i<C_i,\\
    \infty,&T_i\geq C_i.
    \end{cases}
\]
For current-status observations, the recorded interval is $(0,C_i]$ when $T_i<C_i$ and $(C_i,\infty)$ otherwise.

For general interval censoring, a random number of inspection times in $[0,C_i]$ is generated for subjects with $T_i<C_i$. A random intensity $G_i\sim\mathrm{Gamma}(3,1/3)$ is first generated, where the arguments denote shape and rate, followed by $N_i\mid G_i\sim\mathrm{Poisson}(G_i)$. Conditional on $N_i$, inspection times are sampled uniformly on $[0,C_i]$. The observed interval is the smallest inspection interval containing $T_i$. Its left endpoint is the latest inspection before $T_i$, or $0$ if no such inspection exists. Its right endpoint is the earliest inspection after $T_i$, or $C_i$ if no such inspection exists. Subjects with $T_i\geq C_i$ are recorded as $(C_i,\infty)$.

\subsection{Comparison Methods and Tuning}
\label{subsec:simulation-tuning}
Three NRT configurations were compared: a fixed-partition model, an adaptive model with learned cut proportions, and a fixed-partition control using the depth and network architecture selected for the adaptive model. The third configuration assesses the contribution of learning the partition at the same selected architecture. These configurations correspond to the first three rows of each table.

The comparison methods were NNCDE \citep{hu_conditional_2023}, DeepSurv \citep{katzman_deepsurv_2018}, deep distribution regression (DDR) \citep{li_deep_2021}, conditional inference trees (CTree) \citep{hothorn_unbiased_2006,fu_survival_2017}, ICcforest \citep{yao_ensemble_2021}, a linear Cox model, and parametric AFT models. NNCDE and DeepSurv were evaluated for exact and right-censored outcomes, DDR for exact outcomes, and ICcforest for censored outcomes. CTree, Cox, and AFT were evaluated for all observation types shown in the tables.

For each reported setting, a single separately generated data set was used for four-fold cross-validation of NRT, NNCDE, and DeepSurv, where applicable. Validation NLL was used for NRT and NNCDE, and validation concordance was used for DeepSurv. For these three methods, the selected hyperparameters were then held fixed across all $500$ Monte Carlo replications, without repeating cross-validation. For ICcforest, \texttt{mtry} was tuned by out-of-bag evaluation within each replication for every censored-data setting. Detailed tuning strategies, training procedures, and settings for the remaining methods are given in Section~S2 of the Supplementary Material.

\subsection{Simulation and Evaluation Protocol}
\label{subsec:simulation-evaluation}
Each reported setting was evaluated over $500$ Monte Carlo replications generated independently of the tuning data. Each replication contained $5000$ observations and a new random split into $4000$ training and $1000$ test observations. Within a replication, all applicable methods received the same covariates, observed endpoints, and train/test indices.

The primary metric is the average sum of squared errors (ASSE) between the predicted survival function $\hat S$ and the true survival function $S_*$. Let $v_{\max}$ be the largest finite observed endpoint in the training sample. Performance was then evaluated over the test set as
\begin{equation}
  \mathrm{ASSE}
  =
  \frac{1}{n_{\mathrm{test}}}
  \sum_{i=1}^{n_{\mathrm{test}}}
  \sum_{g=1}^{256}
  \{\hat S(t_g\mid x_i)-S_*(t_g\mid x_i)\}^2,
  \label{eq:asse}
\end{equation}
where $n_{\mathrm{test}}=1000$. The true conditional survival function was not supplied as a fitting or validation target; it was used to evaluate ASSE, with the Weibull mixture marginalized over its unobserved component. Each table reports the mean ASSE and the sample standard deviation across the $500$ replications and lower values indicate more accurate survival estimates. 

\subsection{Simulation Results}
The results in Tables~\ref{tab:nrt_exact}--\ref{tab:nrt_current} show that lower mean ASSE was generally obtained with NRT than with the benchmark methods, and learning the response partition often provided a further improvement over fixed-partition NRT. NRT and NRT$^{*}$ use partition-KL and node-mean KL penalties, respectively. Red and blue entries identify the smallest and second-smallest mean ASSE among methods other than NRT$^{*}$, respectively, using their unrounded means; displayed values may coincide after rounding. The discussion below refers to the unstarred results.

Lower mean ASSE was obtained with adaptive NRT than with every applicable benchmark in $27$ of the $29$ reported settings. This pattern was observed across different event-time distributions and observation schemes. In NRT, the conditional distribution is represented through a sequence of tree splits, while nonlinear covariate effects are modeled by the neural network. The comparisons support the use of this construction for conditional survival estimation under the settings considered. Two exceptions were observed. In the PH setting with $q=0.2$ right censoring, mean ASSE was $0.254$ for DeepSurv and $0.273$ for adaptive NRT. For the closer mixture with $q=0.5$ right censoring, lower mean ASSE was obtained with ICcforest ($0.300$), DeepSurv ($0.321$), and NNCDE ($0.426$) than with adaptive NRT ($0.448$).

Adaptive partitioning reduced mean ASSE relative to the independently tuned fixed model in $23$ of the $29$ settings, based on the unrounded means. The comparison with the fixed model at the adaptive architecture provides a more direct assessment of learned cut proportions. In this control, the tree depth and network architecture selected for adaptive NRT were retained, but the cut proportions were fixed at $1/2$ at every internal node. The branching probabilities were still estimated by the neural network. Lower mean ASSE was obtained with adaptive NRT than with this control in $19$ of the $29$ settings. These comparisons support a benefit from learning cut proportions in many settings, although the improvement was not consistent across all observation schemes.

Lower mean ASSE was obtained with adaptive NRT than with both fixed-partition configurations in all eight reported mixture settings. For example, in the uncensored separated mixture, mean ASSE was $0.097$ for adaptive NRT, compared with $0.112$ for the independently tuned fixed model and $0.128$ for the fixed control at the adaptive architecture. Larger reductions were observed for the separated mixture with $q=0.5$ right censoring, where the corresponding values were $0.168$, $0.318$, and $0.345$. Under interval censoring with $q=0.2$, they were $0.104$, $0.198$, and $0.188$. Allowing interval widths to vary with the covariates may help accommodate regions of rapid change in the conditional distribution without requiring the same resolution throughout the response range.

The benefit of adaptive partitioning was less consistent in the PH and current-status settings. In the PH setting with $q=0.5$ right censoring, mean ASSE was $0.646$ for adaptive NRT and $0.615$ for the independently tuned fixed model. Under current-status observation, adaptive NRT improved on the independently tuned fixed model in four of the six settings, but improved on the fixed control at the adaptive architecture only in the log-normal setting with $q=0.2$. Thus, the gains relative to the independently tuned fixed model in these settings were not consistently retained when the selected architecture was held fixed. Several differences between adaptive NRT and the control were small, including values that coincide after rounding. These small differences should be interpreted descriptively rather than as evidence of a clear performance advantage.

The simulations support the use of NRT for conditional survival estimation under the nonlinear covariate effects and observation schemes considered. Learning cut proportions often reduced mean ASSE, with improvements over both fixed-partition configurations in the mixture settings. The results also show that adaptive partitioning did not improve every setting, particularly when compared with the fixed control under current-status observation.
\begin{table}[H]
\centering
\small
\resizebox{\textwidth}{!}{%
\begin{tabular}{lccccc}
\toprule
Method & PH & Log-normal & PO & Multimodal (separate) & Multimodal (closer) \\
\midrule
NRT (fixed NLL CV) & 0.102 (0.049) & 0.086 (0.029) & 0.149 (0.043) & \textcolor{blue}{\textbf{0.112 (0.073)}} & 0.140 (0.098) \\
NRT (adaptive NLL CV) & \textcolor{red}{\textbf{0.085 (0.036)}} & \textcolor{red}{\textbf{0.072 (0.019)}} & \textcolor{blue}{\textbf{0.144 (0.057)}} & \textcolor{red}{\textbf{0.097 (0.055)}} & \textcolor{red}{\textbf{0.084 (0.043)}} \\
NRT (fixed at adaptive architecture) & \textcolor{blue}{\textbf{0.092 (0.038)}} & \textcolor{blue}{\textbf{0.072 (0.023)}} & \textcolor{red}{\textbf{0.133 (0.038)}} & 0.128 (0.048) & \textcolor{blue}{\textbf{0.120 (0.055)}} \\
\addlinespace
CTree & 0.799 (0.170) & 2.675 (0.482) & 1.474 (0.157) & 0.317 (0.023) & 0.435 (0.031) \\
Deep distribution regression & 0.241 (0.082) & 0.254 (0.081) & 0.386 (0.115) & 0.609 (0.162) & 0.310 (0.081) \\
DeepSurv & 0.164 (0.105) & 0.188 (0.043) & 0.349 (0.106) & 0.421 (0.107) & 0.254 (0.043) \\
NNCDE & 0.140 (0.033) & 0.119 (0.026) & 0.151 (0.029) & 0.623 (0.197) & 0.312 (0.045) \\
Cox & 0.822 (0.163) & 2.688 (0.483) & 1.510 (0.109) & 0.347 (0.027) & 0.458 (0.026) \\
AFT & 0.858 (0.168) & 2.689 (0.483) & 1.542 (0.109) & 3.837 (0.103) & 1.005 (0.039) \\
\bottomrule
\end{tabular}
}
\caption{Mean ASSE (standard deviation) over 500 replications for uncensored data, comparing NRT, CTree, deep distribution regression (DDR), DeepSurv, NNCDE, Cox, and AFT models. }
\label{tab:nrt_exact}
\end{table}

\begin{table}[H]
\centering
\small
\resizebox{\textwidth}{!}{%
\begin{tabular}{lccccc}
\toprule
Method & PH & Log-normal & PO & Multimodal (separate) & Multimodal (closer) \\
\midrule
NRT (fixed NLL CV) & 0.316 (0.105) & \textcolor{red}{\textbf{0.124 (0.025)}} & \textcolor{blue}{\textbf{0.221 (0.063)}} & \textcolor{blue}{\textbf{0.123 (0.062)}} & \textcolor{blue}{\textbf{0.151 (0.087)}} \\
NRT (adaptive NLL CV) & \textcolor{blue}{\textbf{0.273 (0.091)}} & \textcolor{blue}{\textbf{0.126 (0.026)}} & \textcolor{red}{\textbf{0.205 (0.040)}} & \textcolor{red}{\textbf{0.094 (0.051)}} & \textcolor{red}{\textbf{0.146 (0.100)}} \\
NRT (fixed at adaptive architecture) & 0.273 (0.090) & 0.135 (0.032) & 0.228 (0.049) & 0.330 (0.036) & 0.173 (0.089) \\
\addlinespace
CTree & 2.360 (0.197) & 4.196 (0.317) & 2.310 (0.174) & 0.325 (0.029) & 0.445 (0.033) \\
ICcforest & 1.250 (0.215) & 2.582 (0.416) & 1.210 (0.207) & 0.209 (0.039) & 0.246 (0.048) \\
DeepSurv & \textcolor{red}{\textbf{0.254 (0.069)}} & 0.381 (0.125) & 0.318 (0.045) & 0.299 (0.042) & 0.242 (0.047) \\
NNCDE & 0.445 (0.076) & 0.208 (0.039) & 0.245 (0.046) & 0.717 (0.213) & 0.325 (0.049) \\
Cox & 2.416 (0.114) & 4.231 (0.281) & 2.350 (0.119) & 0.360 (0.030) & 0.471 (0.028) \\
AFT & 2.476 (0.119) & 4.228 (0.280) & 2.384 (0.117) & 3.683 (0.098) & 1.018 (0.040) \\
\bottomrule
\end{tabular}
}
\caption{Mean ASSE (standard deviation) over 500 replications for right-censored data with $q=0.2$, comparing NRT, CTree, ICcforest, DeepSurv, NNCDE, Cox, and AFT models. }
\label{tab:nrt_right_q0p2}
\end{table}

\begin{table}[H]
\centering
\small
\resizebox{\textwidth}{!}{%
\begin{tabular}{lccccc}
\toprule
Method & PH & Log-normal & PO & Multimodal (separate) & Multimodal (closer) \\
\midrule
NRT (fixed NLL CV) & \textcolor{red}{\textbf{0.615 (0.217)}} & \textcolor{blue}{\textbf{0.480 (0.171)}} & 0.608 (0.179) & 0.318 (0.111) & 0.490 (0.026) \\
NRT (adaptive NLL CV) & 0.646 (0.290) & \textcolor{red}{\textbf{0.427 (0.141)}} & \textcolor{red}{\textbf{0.550 (0.164)}} & \textcolor{red}{\textbf{0.168 (0.078)}} & 0.448 (0.062) \\
NRT (fixed at adaptive architecture) & \textcolor{blue}{\textbf{0.646 (0.273)}} & \textcolor{blue}{\textbf{0.480 (0.171)}} & \textcolor{blue}{\textbf{0.597 (0.183)}} & 0.345 (0.107) & 0.489 (0.026) \\
\addlinespace
CTree & 4.076 (0.268) & 9.672 (0.432) & 5.390 (0.384) & 0.374 (0.050) & 0.479 (0.036) \\
ICcforest & 2.198 (0.407) & 5.526 (0.850) & 2.878 (0.522) & \textcolor{blue}{\textbf{0.316 (0.052)}} & \textcolor{red}{\textbf{0.300 (0.056)}} \\
DeepSurv & 0.818 (0.357) & 0.729 (0.100) & 0.841 (0.355) & 1.305 (0.890) & \textcolor{blue}{\textbf{0.321 (0.056)}} \\
NNCDE & 0.916 (0.187) & 0.547 (0.086) & 0.598 (0.140) & 0.816 (0.187) & 0.426 (0.073) \\
Cox & 4.181 (0.163) & 9.767 (0.283) & 5.494 (0.218) & 0.439 (0.049) & 0.516 (0.036) \\
AFT & 4.197 (0.168) & 9.740 (0.280) & 5.477 (0.215) & 3.251 (0.098) & 1.075 (0.047) \\
\bottomrule
\end{tabular}
}
\caption{Mean ASSE (standard deviation) over 500 replications for right-censored data with $q=0.5$, comparing NRT, CTree, ICcforest, DeepSurv, NNCDE, Cox, and AFT models. }
\label{tab:nrt_right_q0p5}
\end{table}

\begin{table}[H]
\centering
\small
\resizebox{\textwidth}{!}{%
\begin{tabular}{lcccccccc}
\toprule
Method & \multicolumn{5}{c}{$q = 0.2$} & \multicolumn{3}{c}{$q = 0.5$} \\
\cmidrule(lr){2-6} \cmidrule(lr){7-9}
 & PH & Log-normal & PO & Multimodal (separate) & Multimodal (closer) & PH & Log-normal & PO \\
\midrule
NRT (fixed NLL CV) & \textcolor{red}{\textbf{0.309 (0.086)}} & \textcolor{blue}{\textbf{0.152 (0.034)}} & \textcolor{blue}{\textbf{0.262 (0.098)}} & 0.198 (0.030) & \textcolor{blue}{\textbf{0.119 (0.052)}} & \textcolor{red}{\textbf{0.590 (0.158)}} & 0.434 (0.104) & 0.606 (0.268) \\
NRT (adaptive NLL CV) & 0.344 (0.105) & \textcolor{red}{\textbf{0.138 (0.030)}} & \textcolor{red}{\textbf{0.259 (0.098)}} & \textcolor{red}{\textbf{0.104 (0.042)}} & \textcolor{red}{\textbf{0.109 (0.038)}} & \textcolor{blue}{\textbf{0.590 (0.157)}} & \textcolor{red}{\textbf{0.364 (0.095)}} & \textcolor{blue}{\textbf{0.522 (0.182)}} \\
NRT (fixed at adaptive architecture) & \textcolor{blue}{\textbf{0.344 (0.105)}} & 0.170 (0.029) & 0.280 (0.103) & 0.188 (0.048) & 0.270 (0.044) & \textcolor{red}{\textbf{0.590 (0.158)}} & \textcolor{blue}{\textbf{0.365 (0.096)}} & \textcolor{red}{\textbf{0.522 (0.182)}} \\
\addlinespace
CTree & 2.363 (0.187) & 3.880 (0.139) & 2.259 (0.164) & 0.176 (0.018) & 0.186 (0.011) & 4.128 (0.236) & 9.669 (0.531) & 5.409 (0.333) \\
ICcforest & 1.250 (0.217) & 2.420 (0.335) & 1.209 (0.208) & \textcolor{blue}{\textbf{0.137 (0.021)}} & 0.123 (0.019) & 2.191 (0.418) & 5.588 (0.868) & 2.990 (0.500) \\
Cox & 23.868 (0.778) & 36.793 (0.560) & 30.834 (0.434) & 30.444 (0.283) & 40.167 (0.361) & 78.040 (1.295) & 97.863 (1.134) & 87.902 (1.059) \\
AFT & 2.470 (0.111) & 3.892 (0.124) & 2.322 (0.103) & 1.907 (0.030) & 0.402 (0.013) & 4.208 (0.167) & 9.746 (0.297) & 5.457 (0.216) \\
\bottomrule
\end{tabular}
}
\caption{Mean ASSE (standard deviation) over 500 replications for interval-censored data, comparing NRT, CTree, ICcforest, Cox, and AFT models.}
\label{tab:nrt_interval}
\end{table}

\begin{table}[H]
\centering
\small
\resizebox{\textwidth}{!}{%
\begin{tabular}{lcccccc}
\toprule
Method & \multicolumn{3}{c}{$q = 0.2$} & \multicolumn{3}{c}{$q = 0.5$} \\
\cmidrule(lr){2-4} \cmidrule(lr){5-7}
 & PH & Log-normal & PO & PH & Log-normal & PO \\
\midrule
NRT (fixed NLL CV) & \textcolor{red}{\textbf{0.574 (0.198)}} & \textcolor{blue}{\textbf{0.625 (0.182)}} & 0.740 (0.284) & \textcolor{red}{\textbf{0.837 (0.198)}} & 0.929 (0.193) & 0.846 (0.339) \\
NRT (adaptive NLL CV) & 0.607 (0.162) & \textcolor{red}{\textbf{0.494 (0.106)}} & \textcolor{blue}{\textbf{0.591 (0.120)}} & \textcolor{blue}{\textbf{0.838 (0.198)}} & \textcolor{blue}{\textbf{0.817 (0.190)}} & \textcolor{blue}{\textbf{0.672 (0.245)}} \\
NRT (fixed at adaptive architecture) & \textcolor{blue}{\textbf{0.601 (0.158)}} & 0.778 (0.129) & \textcolor{red}{\textbf{0.590 (0.120)}} & \textcolor{red}{\textbf{0.837 (0.198)}} & \textcolor{red}{\textbf{0.816 (0.190)}} & \textcolor{red}{\textbf{0.672 (0.245)}} \\
\addlinespace
CTree & 2.495 (0.176) & 4.006 (0.125) & 2.393 (0.159) & 4.298 (0.263) & 9.870 (0.342) & 5.595 (0.339) \\
ICcforest & 1.441 (0.217) & 2.746 (0.326) & 1.417 (0.209) & 2.472 (0.398) & 6.107 (0.845) & 3.142 (0.513) \\
Cox & 25.158 (1.512) & 36.490 (0.474) & 30.708 (0.412) & 79.241 (1.511) & 97.582 (1.159) & 87.979 (1.019) \\
AFT & 2.529 (0.115) & 3.928 (0.125) & 2.373 (0.124) & 4.202 (0.154) & 9.765 (0.300) & 5.502 (0.219) \\
\bottomrule
\end{tabular}
}
\caption{Mean ASSE (standard deviation) over 500 replications for current-status data, comparing NRT, CTree, ICcforest, Cox, and AFT models. }
\label{tab:nrt_current}
\end{table}

\section{Real World Data Analysis}
\label{sec:rwd}

\section{Discussion}
\label{sec:discuss}

\clearpage

\bibliography{\bibliographyfile}
\end{document}
